\scanpage{50}
\mspage{6}
% archive sheet number: 49
\uline{Remarque} En fait, on a un foncteur plus précis, vers une catégorie de graphes munis d'une structure où la partie $\rho$ de la pondération est \deleted{complétée}\inserted{remplacée} par une structure plus précise, à savoir \deleted{une application}\inserted{une \deleted{\unsure{appl.}} famille}
\[
\deleted{\unsure{A}\ (\widetilde R,\alpha,p:R\longrightarrow\underset{\Vert\atop\vec A/\sigma}{\vec A})}\quad(R_a)_{a\in A}
\]
d'ensembles, et pour $\forall a\in A$, si $\omega(a)$ désigne l'ens. des ``orientations'' de $a$ ($\mathrm{card}(\omega(a))\in\{1,2\}$), une application

% slanted
\margin{$\omega(a)\mathrel{\mathop{=}\limits^{\mathrm{def}}}\{\vec a\in\vec A\mid\vec a\text{ au-dessus de }a\}$}

\deleted{\unsure{$R$ est un ens. fini muni d'une involution de} $R$, et $p$ une application $R\longrightarrow A$ \unsure{compatible à $\alpha$}}
\[
\omega\longmapsto u_\omega:\quad\omega(a)\longrightarrow\mathrm{Circ}(R_a)
\]
ens des permutations circulaires sur $R_a$

\deleted{et induisant une \unsure{bij.} une application \unsure{injective}}
\[
\deleted{R/\alpha\longrightarrow\vec A}
\]
sous la condition que si $\omega,\omega'\in\omega(a)$, $\omega\ne\omega'$, alors $\deleted{u_\omega u_{\omega'}=\mathrm{id}_R}\ u_{\omega'}=u_\omega^{-1}$.

On peut aussi considérer que la donnée des \unsure{suppléments|supplémentaires} revient à la donnée de
\[
\begin{aligned}
r:&\quad A\longrightarrow\mathbb N\\
&\quad a\longmapsto\mathrm{Card}\,p^{-1}(a)
\end{aligned}
\]
(la ``pondération-arêtes''), et la donnée d'une famille de torseurs
\[
(R_a)_{a\in A},\ \deleted{\unsure{\omega\mapsto}}
\]
$R_a$ un torseur sous $\mathbb Z/r(a)\mathbb Z\,\unsure{\mathop{\wedge}\limits_{\{\pm1\}}^{\omega(a)}|\mathop{\wedge}\limits_{\{\pm1\}}\omega(a)}$. On associe à $(X,\widetilde\Sigma,R)$ une telle structure supplémentaire:
\[
R_a=R\cap\widetilde\Sigma_a
\]
($\widetilde\Sigma_a$ la comp. connexe de $\widetilde\Sigma$ \deleted{\unsure{correspondante}}\inserted{qui correspond à $a\in A\simeq\pi_0(\widetilde\Sigma)$})
\[
\deleted{R=\unsure{\widetilde R}(\sigma),\quad R\longrightarrow\underset{\Vert\atop\pi_0(\widetilde\Sigma)}{S}\quad\unsure{\text{la donnée de }R\subset\widetilde\Sigma}}
\]
$\omega\mapsto u_\omega$ défini de façon évidente.

\deleted{et $p$ déduit de l'ordre circulaire sur les $R_a$}

\uline{NB} Même si $r(\omega)=1$ ou même $=2$, le fait de disposer
